%some macro definitions

% format
\newcommand{\class}[1]{{\normalfont\slshape #1\/}}

% entities
\newcommand{\Feup}{Faculdade de Engenharia da Universidade do Porto}

\newcommand{\svg}{\class{SVG}}
\newcommand{\scada}{\class{SCADA}}
\newcommand{\scadadms}{\class{SCADA/DMS}}

%til
%{\raise.17ex\hbox{$\scriptstyle\sim$}}

%piechart
\newcommand{\slice}[4]{
  \pgfmathparse{0.5*#1+0.5*#2}
  \let\midangle\pgfmathresult

  % slice
  \draw[thick,fill=black!10] (0,0) -- (#1:1.0) arc (#1:#2:1.0) -- cycle;

  % outer label
  \node[label=\midangle:#4] at (\midangle:1) {};

  % inner label
  \pgfmathparse{min((#2-#1-10.0)/110.0*(-0.3),0)}
  \let\temp\pgfmathresult
  \pgfmathparse{max(\temp,-0.5) + 0.8}
  \let\innerpos\pgfmathresult
  \node at (\midangle:\innerpos) {#3};
}



\newcommand{\rowgroup}[1]{\hspace{-2em}#1}  % for indentation
\newcolumntype{C}{>{\centering\arraybackslash}X}
\newcommand\mc[1]{\multicolumn{1}{C}{#1}}
%\sisetup{group-separator={,} , group-four-digits=true}

\newcommand*\rfrac[2]{{}^{#1}\!/_{#2}}